% Part: normal-modal-logic
% Chapter: completeness
% Section: lindenbaums-lemma

\documentclass[../../../include/open-logic-section]{subfiles}

\begin{document}

\olfileid[hi]{nml}{com}{lin}

\olsection{लिंडेनबाउम का लेम्मा}

लिंडेनबाउम का लेम्मा स्थापित करता है कि !!{formula}ों का प्रत्येक
$\Sigma$-संगत समुच्चय कम-से-कम एक \emph{पूर्ण} $\Sigma$-संगत समुच्चय में
समाहित होता है। विहित प्रतिरूप का हमारा निर्माण दिखाएगा कि प्रत्येक पूर्ण
$\Sigma$-संगत समुच्चय~$\Delta$ के लिए विहित प्रतिरूप में ऐसा जगत है जहाँ
ठीक $\Delta$ के सभी !!{formula} सत्य हैं। अतः लिंडेनबाउम का लेम्मा प्रत्याभूत
करता है कि प्रत्येक $\Sigma$-संगत समुच्चय विहित प्रतिरूप के किसी जगत पर
सत्य है।

\begin{thm}[लिंडेनबाउम का लेम्मा]\ollabel{thm:lindenbaum}
  यदि $\Gamma$, $\Sigma$-संगत है, तो $\Gamma$ को विस्तृत करने वाला कोई
  पूर्ण $\Sigma$-संगत समुच्चय $\Delta$ है।
\end{thm}

\begin{proof}
मान लें कि $!A_0$, $!A_1$, \dots{} भाषा के सभी सूत्रों की समग्र सूची है
(पुनरावृत्तियाँ अनुमत हैं)। उदाहरणार्थ, सूची को $\Obj p_0$ से आरंभ करें,
और प्रत्येक चरण $n \ge 1$ पर केवल $\Obj p_0$, \dots,~$\Obj p_n$ में से
चरों का उपयोग करने वाले लंबाई~$n$ के परिमिततः अनेक सूत्र सूचीबद्ध करें।
हम $n$ पर आगमन से !!{formula}ों के समुच्चय $\Delta_n$ परिभाषित करते हैं और
फिर $\Delta = \bigcup_n \Delta_n$ रखते हैं। पहले $\Delta_0 = \Gamma$
रखें। मान लें कि $\Delta_n$ परिभाषित हो चुका है; तब $\Delta_{n+1}$ को इस
प्रकार परिभाषित करते हैं:
\[
\Delta_{n+1} =
\begin{cases}
  \Delta_n \cup \{!A_n\}, & \text{यदि $\Delta_n \cup \{ !A_n\}$
    $\Sigma$-संगत है;} \\
  \Delta_n \cup \{ \lnot !A_n\}, & \text{अन्यथा।}
\end{cases}
\]
अब $\Delta = \bigcup_{n=0}^\infty \Delta_n$ रखें।

हमें दिखाना है कि यह परिभाषा वास्तव में अपेक्षित गुणों वाला समुच्चय
$\Delta$ देती है, अर्थात् $\Gamma \subseteq \Delta$ और $\Delta$ पूर्ण
$\Sigma$-संगत है।

स्पष्ट है कि $\Gamma \subseteq \Delta$, क्योंकि निर्माण से $\Delta_0
\subseteq \Delta$, और $\Delta_0 = \Gamma$। वास्तव में प्रत्येक~$n$ के
लिए $\Delta_n \subseteq \Delta$, क्योंकि $\Delta$ सभी~$\Delta_n$ का
सम्मिलन है। (निर्माण के प्रत्येक चरण में हम पहले से बने समुच्चय में
!!a{formula} जोड़ते हैं, इसलिए $\Delta_n \subseteq \Delta_{n+1}$; और
$\subseteq$ के संक्रमणशील होने से जब भी $n \le m$, तब $\Delta_n
\subseteq \Delta_m$।) निर्माण के प्रत्येक चरण में हम $!A_n$ या $\lnot
!A_n$ में से एक जोड़ते हैं, और प्रत्येक !!{formula} सभी~$!A_n$ की सूची में
(कम-से-कम एक बार) आता है। अतः प्रत्येक $!A$ के लिए या तो $!A \in
\Delta$, या $\lnot !A \in \Delta$; इसलिए परिभाषा से $\Delta$ पूर्ण है।

अंततः हमें दिखाना है कि $\Delta$, $\Sigma$-संगत है। इसके लिए हम दिखाते
हैं कि (क) यदि $\Delta$, $\Sigma$-असंगत होता, तो कोई $\Delta_n$ भी
$\Sigma$-असंगत होता, और (ख) सभी $\Delta_n$, $\Sigma$-संगत हैं।

मान लें कि $\Delta$, $\Sigma$-असंगत है। तब $\Delta \Proves[\Sigma]
\lfalse$, अर्थात् ऐसे $!A_1$, \dots,~$!A_k \in \Delta$ हैं कि $\Sigma
\Proves !A_1 \lif (!A_2 \lif \cdots (!A_k \lif \lfalse)\dots)$। चूँकि
$\Delta = \bigcup_{n=0}^\infty \Delta_n$, प्रत्येक $!A_i$ किसी~$n_i$ के
लिए $\Delta_{n_i}$ में है। इन $n_i$ में सबसे बड़े को $n$ मानें। चूँकि
$n_i \le n$, इसलिए $\Delta_{n_i} \subseteq \Delta_n$। अतः सभी $!A_i$
किसी एक~$\Delta_n$ में हैं। इसका अर्थ होगा $\Delta_n \Proves[\Sigma]
\lfalse$, अर्थात् $\Delta_n$, $\Sigma$-असंगत है।

प्रत्येक $\Delta_n$ की $\Sigma$-संगति दिखाने के लिए हम $n$ पर सरल आगमन
का उपयोग करते हैं। $\Delta_0 = \Gamma$, और हमने माना था कि $\Gamma$,
$\Sigma$-संगत है। अतः दावा $n = 0$ के लिए सत्य है। अब मान लें कि वह $n$
के लिए सत्य है, अर्थात् $\Delta_n$, $\Sigma$-संगत है। यदि $\Delta_n \cup
\{!A_n\}$, $\Sigma$-संगत है तो $\Delta_{n+1}$ यही है; अन्यथा वह
$\Delta_n \cup \{\lnot!A_n\}$ है। पहली स्थिति में $\Delta_{n+1}$
स्पष्टतः $\Sigma$-संगत है। किंतु
\olref[prf][con]{prop:consistencyfacts}\olref[prf][con]{prop:consistencyfacts-c}
के अनुसार $\Delta_n \cup \{!A_n\}$ और $\Delta_n \cup \{\lnot!A_n\}$
में से कोई एक संगत है, इसलिए दूसरी स्थिति में भी $\Delta_{n+1}$ संगत है।
\end{proof}

\begin{cor}\ollabel{cor:provability-characterization}
  $\Gamma \Proves[\Sigma] !A$ तभी जब $\Gamma$ को विस्तृत करने वाले
  प्रत्येक पूर्ण $\Sigma$-संगत समुच्चय $\Delta$ के लिए $!A \in \Delta$
  (इसमें $\Gamma = \emptyset$ की स्थिति भी है, जिसमें हमें मोडल तंत्र
  $\Sigma$ का एक और अभिलक्षण मिलता है)।
\end{cor}

\begin{proof}
  मान लें $\Gamma \Proves[\Sigma] !A$, और $\Delta$, $\Gamma$ को विस्तृत
  करने वाला कोई पूर्ण $\Sigma$-संगत समुच्चय हो। यदि $!A \notin \Delta$,
  तो अधिकतमता से $\lnot!A \in \Delta$; और इसलिए $\Delta
  \Proves[\Sigma] !A$ (एकदिष्टता से) तथा $\Delta \Proves[\Sigma]
  \lnot!A$ (स्वपरावर्तिता से), अतः $\Delta$ असंगत है। विलोमतः, यदि
  $\Gamma \Proves/[\Sigma] !A$, तो $\Gamma \cup \{ \lnot!A\}$,
  $\Sigma$-संगत है; और लिंडेनबाउम के लेम्मा से $\Gamma \cup \{ \lnot!A
  \}$ को विस्तृत करने वाला कोई पूर्ण संगत समुच्चय $\Delta$ है। संगति से
  $!A \notin \Delta$।
\end{proof}

\end{document}
